% Generated draft: requires mathematical and source review.
\documentclass[11pt]{article}
\usepackage[margin=25mm]{geometry}
\usepackage{fontspec}
\setmainfont{Latin Modern Roman}
\usepackage{amsmath,amssymb,mathtools,unicode-math}
\setmathfont{Latin Modern Math}
\usepackage{enumitem,needspace,xurl,hyperref}
\setlength{\parindent}{0pt}
\setlength{\parskip}{5pt}
\setlength{\emergencystretch}{3em}
\newcommand{\N}{\mathbb{N}}
\newcommand{\Z}{\mathbb{Z}}
\newcommand{\Q}{\mathbb{Q}}
\newcommand{\R}{\mathbb{R}}
\newcommand{\C}{\mathbb{C}}
\newcommand{\F}{\mathbb{F}}
\newcommand{\A}{\mathbb{A}}
\newcommand{\PP}{\mathbb{P}}
\newcommand{\E}{\mathbb{E}}
\newcommand{\eps}{\varepsilon}
\newcommand{\dd}{\,\mathrm{d}}
\begin{document}
\section*{UnsolvedMath ID: 30001446}
\textbf{Source problem code:} catalogue entry
\subsection{Definitions and mathematical statement}
A partition \(\lambda=(\lambda_1\geq\cdots\geq\lambda_\ell>0)\) has size \(|\lambda|=\sum_i\lambda_i\). Let \(X=(x_1,x_2,\ldots)\) be a formal alphabet, let \(q\) be an indeterminate, and write \(s_\mu(X)\) for the Schur symmetric function indexed by a partition \(\mu\). For a homogeneous symmetric function \(F\) of degree \(n\) over \(\Q(q)\), define its Schur coefficients by
\[
F=\sum_{\mu\vdash n}a_\mu(q)s_\mu(X),\qquad [s_\mu]F=a_\mu(q).
\]
Here \(\N[q]\) means polynomials in \(q\) with nonnegative integer coefficients.

The supplied question is: for every partition \(\lambda\vdash n\) and every \(k\in\N\), does
\[
[s_\lambda]\!\left(\frac{J_\lambda(X;q,q^k)}{(1-q)^n}\right)\in\N[q]
\]
hold? The same partition \(\lambda\) indexes both \(J_\lambda\) and the Schur coefficient being requested.

The symbol \(J_\lambda(X;q,t)\) denotes the Macdonald symmetric function referred to in the supplied record, but that record does not specify its normalization or any alphabet convention beyond the displayed formula. Thus the question above is conditional on the intended meaning of \(J_\lambda\) and on the specialization being defined. The record also leaves unspecified whether \(0\in\N\) for the parameter \(k\), and whether the empty partition is included. Its dated background assessment says that the actual conjectural statement is missing, whereas its statement-verification field reports that the displayed question was checked against the official report. These conflicting source qualifications remain unresolved.
\subsection{Short English statement}
For each partition \(\lambda\) of \(n\) and natural number \(k\), the supplied question asks whether the coefficient of \(s_\lambda\) in \(J_\lambda(X;q,q^k)/(1-q)^n\) is a polynomial with nonnegative integer coefficients. The normalization and parameter boundary conventions require clarification, and the supplied source assessments conflict about the completeness of this formulation.
\subsection{Sources}
\begin{itemize}
\item Supplied local registry record, permanent UnsolvedMath ID \(30001446\), exact source problem code \texttt{catalogue entry}, title \emph{Positivity in Specialized Macdonald-Polynomial Expansions}.
\item \emph{Combinatorial Representation Theory} (2010), Oberwolfach Reports, pp.~799--882, DOI \texttt{10.4171/owr/2010/15}. Supplied source URL and literature-source link labelled ``original OWR report'': \url{https://doi.org/10.4171/owr/2010/15}.
\item The supplied background associates the fragment with Ram's positivity conjecture for Macdonald-polynomial expansions. The proposer field is empty; the record supplies no fuller attribution or pinpoint location within the report.
\item Literature assessment dated exactly \texttt{2026-08-21}: the supplied curation labels the problem open and says that the extracted fragment omits the actual conjectural statement, preventing verification of a complete resolved target. This is supplied dated curation, not an independently verified literature assessment.
\item The record separately gives statement status \texttt{corrected\_verified} and reports that its mathematical content was checked against the official report, with notation, context, and question form normalized. This verification claim is supplied metadata and has not been independently checked here.
\end{itemize}
\end{document}
